% ============================================================
% M2 轮4 成卷·S6 件2 —— 滚动测评卷（A）（16 题＝命制-滚A 全卷题面入卷）
% 2026-09-12｜执行：S6 成卷代理（Kimi K2.8 第三臂首派，看板九十四）
% 样式：卷式宏承 v11 测评卷 qp-* 六宏（原样拷贝、零改动）＋练习线题块；走 L1 参数组合登记（R2 口径）。
% 题源：命制/命制-滚A-16题.md 全卷 16 题（复核-滚A-glm 15 通过＋滚A-11 主脑翻案通过；0 键离）；
%   主脑裁决义务随卷落实 2 处：滚A-12 首空「一定不＿＿」显式化为「一定不平行」；滚A-16(2)「若有 t」
%   措辞改「判断∠ABP 与 t 的关系并指出直角顶点」；滚A-6↔W7 不同卷（W7 不在本卷，天然满足）。
% 卷式（照命制件头注）：单选 7×5＝35｜多选 2×5＝10｜填空 3×5＝15｜解答 8＋10＋10＋12＝40；40 分钟 100 分。
% 全卷无图（三维图禁绘，命制件图占位默认转「无图」，题面文字自足）；答案制A＝题答分离，卷末答案速查表。
% ============================================================
\documentclass[fontset=none]{ctexart}
\usepackage{qp-m3}
\ansblockgrayfalse
%% 修因探针：六宏 qp-layout.tex:130 全局 extrarowheight=1.7mm，qp-m3 未复刻
\setlength{\extrarowheight}{1.7mm}

% ------------------------------------------------------------
% 件型层 A｜页面：8 开横放三栏（承测评卷件型层）
% ------------------------------------------------------------
\geometry{paperwidth=399.8mm,paperheight=284.2mm,left=8.9mm,right=8.9mm,
  top=12.0mm,bottom=17.6mm,twoside,headheight=0pt,headsep=0pt,footskip=7.1mm}
\setlength{\topskip}{0pt}
\raggedbottom
\renewcommand{\normalsize}{\fontsize{10.5pt}{17.7pt}\selectfont}
\linespread{1}\normalsize
\setlength{\parindent}{0pt}
\setlength{\parskip}{0pt}
\newlength{\jpcolw}\setlength{\jpcolw}{120.0mm}
\newlength{\jpgap}\setlength{\jpgap}{11.0mm}
\xeCJKsetup{CJKglue={\hskip 0.04em plus 0.04em minus 0.01em}}

\definecolor{silklight}{gray}{0.8235}
\definecolor{silkdark}{gray}{0.6510}
\definecolor{juanrule}{gray}{0.1255}
\definecolor{subgrey}{gray}{0.4}

% ------------------------------------------------------------
% 件型层 B｜JT-01 丝带角饰（照抄测评卷件型层）
% ------------------------------------------------------------
\newcommand{\juanmathSize}{17pt}
\newcommand{\sishou}{%
  \begin{tikzpicture}[overlay,remember picture,x=1mm,y=1mm]
    \path (current page.north west) coordinate (O);
    \begin{scope}[shift={(O)}]
      \fill[silklight] (6.9,-9.39) -- (40.1,-9.39) -- (6.9,-42.59) -- cycle;
      \draw[white,line width=0.28mm] (9.21,-40.29) -- (9.21,-11.63) -- (37.87,-11.63);
      \fill[silkdark] (22.94,-8.21) -- (34.15,-8.21) -- (5.78,-36.58) -- (5.78,-25.37) -- cycle;
      \fill[silkdark] (34.15,-8.21) -- (34.15,-9.20) -- (33.16,-9.20) -- cycle;
      \fill[silkdark] (5.78,-36.58) -- (6.78,-36.58) -- (6.78,-35.58) -- cycle;
      \node[white,inner sep=0pt,anchor=center,rotate=45] at (17.23,-19.74)
        {\fontsize{\juanmathSize}{\juanmathSize}\selectfont\heitiao 数学};
    \end{scope}
  \end{tikzpicture}}

% ------------------------------------------------------------
% 件型层 C｜卷头零件＋题区零件（承测评卷件型层；卷名/副题/时间/组行按滚A 卷式改文）
% ------------------------------------------------------------
\newcommand{\zeroheight}[1]{\raisebox{0pt}[0pt][0pt]{#1}}
\newcommand{\jpcen}[1]{\hbox to\jpcolw{\hfil\zeroheight{#1}\hfil}\nointerlineskip}
\newcommand{\jpright}[1]{\hbox to\jpcolw{\hfill\zeroheight{#1}}\nointerlineskip}
\newcommand{\vgt}[1]{\par\vskip\dimexpr#1-6.25mm\relax}

\newcommand{\tianxieg}[1]{{\fontsize{\qpJTbLabelSize}{13pt}\selectfont\heitiao #1：}%
  \rule[-0.62mm]{\qpJTbLineLen}{0.2mm}}
\newcommand{\jtianxie}{\tianxieg{班级}\hspace{\qpJTbGapGrp}\tianxieg{姓名}%
  \hspace{\qpJTbGapGrp}\tianxieg{得分}}
\newcommand{\juanmingSize}{20.5pt}
\newcommand{\jjuanming}{{\fontsize{\juanmingSize}{26pt}\selectfont\heizhang
  \renewcommand{\CJKglue}{\hskip 0pt}滚动测评卷（A）}}
\newcommand{\jjunrule}{{\color{juanrule}\rule{58.0mm}{\qpJTcRule}}}
\newcommand{\jfuti}{{\fontsize{\qpJTdSize}{18pt}\selectfont\heitiao\color{subgrey}%
  \renewcommand{\CJKglue}{\hskip 0pt}第一章（1.1.1–1.1.3）}}
\newcommand{\jptime}{{\fontsize{\qpJTeSize}{17.7pt}\selectfont（时间：40分钟\quad 分值：100分）}}
\newcommand{\zuhang}[2]{\par\noindent\hangindent=5.7mm\hangafter=1
  {\fontsize{\qpJTfSize}{17.7pt}\selectfont\heitiao #1}#2\par}
\newcommand{\ti}[2]{\par\noindent\hangindent=5.7mm\hangafter=1
  {\fontsize{11.4pt}{17.7pt}\selectfont\heihao{\numboldjian #1}.}\hspace{2.7mm}#2\par}
\newcommand{\fenzhi}[1]{（#1分）}
\newcommand{\optline}[1]{\par\noindent\hangindent=5.7mm\hangafter=0
  {\fontsize{10.5pt}{19pt}\selectfont #1}\par}
\newcommand{\optII}[2]{\makebox[53.05mm][l]{#1}#2}
\newcommand{\optIV}[4]{\makebox[21.5mm][l]{#1}\makebox[21.5mm][l]{#2}\makebox[21.5mm][l]{#3}#4}
\newcommand{\kw}{\hfill（\hspace{3.6mm}）\hspace*{5.4mm}}
\newcommand{\qpart}[1]{\par\noindent\hangindent=5.7mm\hangafter=0
  \hspace*{5.7mm}{\fontsize{10.5pt}{17.7pt}\selectfont #1}\par}

\newdimen\hA \hA=3.03mm
\newdimen\hB \hB=12.12mm
\newdimen\hC \hC=1.89mm
\newdimen\hD \hD=5.95mm
\newdimen\hE \hE=8.04mm
\newdimen\hF \hF=5.70mm
\newdimen\hG \hG=6.35mm
\newcommand{\jphead}{\hbox to\jpcolw{\sishou\hfill}\nointerlineskip
  \vskip\hA\jpright{\jtianxie}%
  \vskip\hB\jpcen{\jjuanming}%
  \vskip\hC\jpcen{\jjunrule}%
  \vskip\hD\jpcen{\jfuti}%
  \vskip\hE\jpcen{\jptime}%
  \vskip\hF
  \zuhang{一、选择题}{：本题共7小题，每小题5分，共35分．\ 在每小题给出的四个选项中，只有一项是符合题目要求的．}%
  \vgt{\hG}}

% ------------------------------------------------------------
% 件型层 D｜YJ-02 文字行页脚（同测评卷；卷名改滚A）
% ------------------------------------------------------------
\newcommand{\jpnum}[1]{{\fontsize{\qpYJbNumSize}{\qpYJbNumSize}\selectfont\numbold #1}}
\newcommand{\jptxtsize}{7.5pt}
\newcommand{\jptxt}{\fontsize{\jptxtsize}{\jptxtsize}\selectfont}
\newcommand{\juanname}{滚动测评卷（A）}
\newcommand{\pqt}{\ifnum\value{page}<10 0\fi\thepage}
\setlength{\headwidth}{\textwidth}
\fancyfoot[RO]{\jptxt\juanname\hspace{3.95mm}{\heitiao 滚动卷}\hspace{3.88mm}卷\hspace{2.25mm}\jpnum{\pqt}}
\fancyfoot[LE]{\jptxt\jpnum{\pqt}\hspace{1.9mm}卷\hspace{3.95mm}{\heitiao 羿郭工作室}%
  \hspace{3.0mm}高中数学\hspace{2.75mm}选择性必修第一册\hspace{2.8mm}RJB}

% ------------------------------------------------------------
% 件型层 E｜三栏排布器＋卷末答案速查表
% ------------------------------------------------------------
\newcommand{\jpcol}[1]{\raisebox{0pt}[0pt][0pt]{\vtop{\hsize\jpcolw\parskip0pt\parindent0pt
  \setlength\linewidth{\jpcolw}\setlength\columnwidth{\jpcolw}   % 换装补钉：qp-m3 括线盒按 \linewidth 取宽
  \fontsize{10.5pt}{17.7pt}\selectfont\relax#1}}}
\newcommand{\jpthree}[3]{\nointerlineskip\noindent\jpcol{#1}\hspace{\jpgap}\jpcol{#2}%
  \hspace{\jpgap}\jpcol{#3}\par}
\newcommand{\jplead}{\hbox{}\nointerlineskip}

\newcommand{\datu}[1]{{\fontsize{9pt}{12pt}\selectfont #1}}
\newcommand{\dabiao}{%
  \par\vgt{\hG}\noindent{\fontsize{10.5pt}{17.7pt}\selectfont\heitiao 答案速查}\par\nopagebreak
  \vskip 1.2mm\noindent\datu{\setlength{\tabcolsep}{1.4pt}%
  \begin{tabular}{|c|*{16}{c|}}
  \hline
  题号 & 1 & 2 & 3 & 4 & 5 & 6 & 7 & 8 & 9 & 10 & 11 & 12 & 13 & 14 & 15 & 16 \\ \hline
  答案 & A & A & D & A & D & A & D & ABC & ABC & \((3,0,0)\) & \(\frac{1}{6}\) & 略 & 略 & 略 & 略 & 略 \\ \hline
  \end{tabular}}\par}

% —— 修因探针2：件面末位再注册 \parallel 的 TikZ 复刻（qp-m3:105 的注册被 unicode-math 覆盖）
\AtBeginDocument{\renewcommand{\parallel}{\mathrel{\begin{tikzpicture}[x=1em,y=1em,baseline={(0,0)}]
  \fill (-0.412,-0.016) -- (-0.375,-0.016) -- (0.114,0.956) -- (0.077,0.956) -- cycle;
  \fill (-0.115,-0.016) -- (-0.078,-0.016) -- (0.411,0.956) -- (0.374,0.956) -- cycle;
\end{tikzpicture}}}}
\begin{document}
% ===================== 第 1 页：栏1 卷头＋Q1–Q3；栏2 Q4–Q7；栏3 组行二＋Q8–Q9＋组行三＋Q10 =====================
\jpthree{\jphead
  \ti{1}{设 \(\overrightarrow{a}\)，\(\overrightarrow{b}\) 为非零向量，则「\(\left|\overrightarrow{a}+\overrightarrow{b}\right|=\left|\overrightarrow{a}\right|+\left|\overrightarrow{b}\right|\)」是「\(\overrightarrow{a}\) 与 \(\overrightarrow{b}\) 共线」的\kw}
  \optline{\optII{A．充分不必要条件；}{B．必要不充分条件；}}
  \optline{\optII{C．充要条件；}{D．既不充分也不必要条件．}}
  \ti{2}{已知两个互相垂直的共点力 \(\overrightarrow{F_1}\)，\(\overrightarrow{F_2}\)，\(\left|\overrightarrow{F_1}\right|=3\,\mathrm{N}\)，\(\left|\overrightarrow{F_2}\right|=4\,\mathrm{N}\)，它们的合力为 \(\overrightarrow{F}=\overrightarrow{F_1}+\overrightarrow{F_2}\)，则合力 \(\overrightarrow{F}\) 与 \(\overrightarrow{F_1}\) 夹角的余弦值为\kw}
  \optline{\optIV{A．\(\frac{3}{5}\)；}{B．\(\frac{4}{5}\)；}{C．\(\frac{3}{4}\)；}{D．\(\frac{4}{3}\)．}}
  \ti{3}{在四面体 \(A-BCD\) 中，设 \(\overrightarrow{AB}=\overrightarrow{a}\)，\(\overrightarrow{AC}=\overrightarrow{b}\)，\(\overrightarrow{AD}=\overrightarrow{c}\)，则下列向量中能与 \(\overrightarrow{a}\)，\(\overrightarrow{b}\) 构成空间一个基底的是\kw}
  \optline{\optII{A．\(\overrightarrow{BD}+\overrightarrow{DA}\)；}{B．\(\overrightarrow{AB}-\overrightarrow{AC}\)；}}
  \optline{\optII{C．\(\overrightarrow{BC}+\overrightarrow{CD}+\overrightarrow{DA}+\overrightarrow{AB}\)；}{D．\(\overrightarrow{BC}+\overrightarrow{BD}\)．}}
}{\jplead
  \ti{4}{空间直角坐标系中，\(A(1,1,0)\)，\(B(1,-1,0)\)，\(C(0,0,\sqrt{2})\)，则 \(\triangle ABC\) 的形状是\kw}
  \optline{\optII{A．等边三角形；}{B．直角三角形；}}
  \optline{\optII{C．等腰非等边三角形；}{D．既非等腰也非直角三角形．}}
  \ti{5}{空间直角坐标系中，点 \(P(m-1,\ m+2,\ 0)\)（\(m\in\mathbf{R}\)），则 \(P\) 一定不可能\kw}
  \optline{\optII{A．在 \(x\) 轴上；}{B．在 \(y\) 轴上；}}
  \optline{\optII{C．在坐标平面 \(xOy\) 内；}{D．为坐标原点．}}
  \ti{6}{空间四边形 \(ABCD\) 中，\(E\)，\(F\) 分别为对角线 \(AC\)，\(BD\) 的中点，设 \(\overrightarrow{AB}=\overrightarrow{a}\)，\(\overrightarrow{BC}=\overrightarrow{b}\)，\(\overrightarrow{CD}=\overrightarrow{c}\)，\(\overrightarrow{DA}=\overrightarrow{d}\)，则 \(\overrightarrow{EF}=\)\kw}
  \optline{\optII{A．\(\frac{1}{2}(\overrightarrow{a}+\overrightarrow{c})\)；}{B．\(\frac{1}{2}(\overrightarrow{b}+\overrightarrow{d})\)；}}
  \optline{\optII{C．\(\frac{1}{4}(\overrightarrow{a}+\overrightarrow{b}+\overrightarrow{c}+\overrightarrow{d})\)；}{D．\(\frac{1}{2}(\overrightarrow{a}-\overrightarrow{c})\)．}}
  \ti{7}{设 \(\overrightarrow{a}\)，\(\overrightarrow{b}\)，\(\overrightarrow{c}\) 均为非零向量，下列由平面向量类比到空间向量的结论中，错误的是\kw}
  \optline{\optII{A．若 \(\overrightarrow{a}\parallel\overrightarrow{b}\)，\(\overrightarrow{b}\parallel\overrightarrow{c}\)，则 \(\overrightarrow{a}\parallel\overrightarrow{c}\)；}{B．若 \(\overrightarrow{a}=\overrightarrow{b}\)，\(\overrightarrow{b}=\overrightarrow{c}\)，则 \(\overrightarrow{a}=\overrightarrow{c}\)；}}
  \optline{\optII{C．\(\left|\overrightarrow{a}\cdot\overrightarrow{b}\right|\leq\left|\overrightarrow{a}\right|\left|\overrightarrow{b}\right|\)；}{D．若 \(\overrightarrow{a}\perp\overrightarrow{b}\)，\(\overrightarrow{b}\perp\overrightarrow{c}\)，则 \(\overrightarrow{a}\parallel\overrightarrow{c}\)．}}
}{\jplead
  \zuhang{二、选择题}{：本题共2小题，每小题5分，共10分．\ 在每小题给出的选项中，有多项符合题目要求．（全部选对得5分，部分选对得2分，有选错的得0分）}%
  \vgt{\hG}
  \ti{8}{空间直角坐标系中，已知 \(A(0,0,0)\)，\(B(2,-2,1)\)，\(C(1,2,2)\)，\(D(4,m,5)\)（\(m\in\mathbf{R}\)），则下列说法正确的是\kw}
  \optline{①\(\overrightarrow{AB}\perp\overrightarrow{AC}\)；②\(\left|\overrightarrow{AB}\right|=\left|\overrightarrow{AC}\right|\)；}
  \optline{③若 \(m=2\)，则 \(A\)，\(B\)，\(C\)，\(D\) 四点共面；④若 \(m=2\)，则 \(\overrightarrow{AD}\parallel\overrightarrow{BC}\)．}
  \optline{\optIV{A．①；}{B．②；}{C．③；}{D．④．}}
  \ti{9}{设 \(P\) 为 \(\triangle ABC\) 所在平面内一点（\(P\) 与三顶点均不重合），则下列说法正确的是\kw}
  \optline{A．若 \(\overrightarrow{PA}+\overrightarrow{PB}+\overrightarrow{PC}=\overrightarrow{0}\)，则 \(P\) 是 \(\triangle ABC\) 的重心；}
  \optline{B．若 \(\overrightarrow{PA}\cdot\overrightarrow{PB}=\overrightarrow{PB}\cdot\overrightarrow{PC}=\overrightarrow{PC}\cdot\overrightarrow{PA}\)，则 \(P\) 是 \(\triangle ABC\) 的垂心；}
  \optline{C．若 \(\left|\overrightarrow{PA}\right|=\left|\overrightarrow{PB}\right|=\left|\overrightarrow{PC}\right|\)，则 \(P\) 是 \(\triangle ABC\) 的外心；}
  \optline{D．若 \(O\) 为平面内任一点，\(\overrightarrow{OP}=\frac{1}{3}(\overrightarrow{OA}+\overrightarrow{OB}+\overrightarrow{OC})\)，则 \(P\) 是 \(\triangle ABC\) 的内心．}
  \vgt{\hG}
  \zuhang{三、填空题}{：本题共3小题，每小题5分，共15分．}%
  \vgt{\hG}
  \ti{10}{在 \(x\) 轴上求一点 \(M\)，使 \(M\) 到 \(A(1,2,3)\) 与 \(B(3,4,1)\) 的距离相等，则 \(M\) 的坐标为\kongbai}
}
\newpage
% ===================== 第 2 页：栏1 Q11–Q12＋组行四；栏2 Q13–Q14；栏3 Q15–Q16＋答案速查表 =====================
\jpthree{\jplead
  \ti{11}{空间中 \(O\)，\(A\)，\(B\)，\(C\) 四点不共面，\(P\) 为空间一点，且 \(\overrightarrow{OP}=\frac{1}{3}\overrightarrow{OA}+\frac{1}{2}\overrightarrow{OB}+t\overrightarrow{OC}\)（\(t\) 为实数）．若 \(P\)，\(A\)，\(B\)，\(C\) 四点共面，则 \(t=\)\kongbai}
  \ti{12}{边长为 \(2\) 的等边三角形 \(ABC\) 中，\(D\) 为 \(BC\) 的中点．将 \(\triangle ABD\) 沿 \(AD\) 折起到 \(\triangle AB'D\) 的位置，折起后记 \(\angle B'DC=\varphi\)（\(0^\circ<\varphi<180^\circ\)），得空间四面体骨架图形 \(A-B'CD\)．则：}
  \qpart{（1）\(\overrightarrow{AB'}\) 与 \(\overrightarrow{DC}\) 的位置关系：两直线一定不平行；\(\overrightarrow{AB'}\cdot\overrightarrow{DC}=\)\kongbai（用含 \(\varphi\) 的式子表示）；当 \(\overrightarrow{AB'}\perp\overrightarrow{DC}\) 时，\(\varphi=\)\kongbai；}
  \qpart{（2）在 \(\overrightarrow{AB'}\perp\overrightarrow{DC}\) 的条件下，\(\left|\overrightarrow{CB'}\right|=\)\kongbai；\(\cos\langle\overrightarrow{AB'},\overrightarrow{CB'}\rangle=\)\kongbai．}
  \vgt{\hG}
  \zuhang{四、解答题}{：本题共4小题，共40分．\ 解答应写出文字说明、证明过程或演算步骤．}%
  \vgt{\hG}
}{\jplead
  \ti{13}{\fenzhi{8}用向量法证明：平行四边形两条对角线的平方和等于它四条边的平方和．（设平行四边形为 \(ABCD\)，\(AC\)，\(BD\) 为两条对角线，四边为 \(AB\)，\(BC\)，\(CD\)，\(DA\)．）}
  \ti{14}{\fenzhi{10}四面体 \(ABCD\) 中，设 \(\overrightarrow{AB}=\overrightarrow{a}\)，\(\overrightarrow{AC}=\overrightarrow{b}\)，\(\overrightarrow{AD}=\overrightarrow{c}\)．\(E\) 为 \(CD\) 的中点，点 \(P\) 在线段 \(BE\) 上，且 \(BP=2PE\)．}
  \qpart{（1）用 \(\overrightarrow{a}\)，\(\overrightarrow{b}\)，\(\overrightarrow{c}\) 表示 \(\overrightarrow{AP}\)；}
  \qpart{（2）若点 \(Q\) 满足 \(\overrightarrow{AQ}=\frac{2}{3}\overrightarrow{a}+\frac{1}{3}\overrightarrow{b}+\frac{1}{3}\overrightarrow{c}\)，判断 \(A\)，\(P\)，\(Q\) 三点是否共线，并说明理由．}
}{\jplead
  \ti{15}{\fenzhi{10}棱长为 \(2\) 的正方体 \(ABCD-A_1B_1C_1D_1\) 中，以 \(D\) 为原点，\(DA\)，\(DC\)，\(DD_1\) 所在直线分别为 \(x\) 轴、\(y\) 轴、\(z\) 轴，建立空间直角坐标系．}
  \qpart{（1）写出 \(A\)，\(C\)，\(B_1\)，\(D_1\) 四点的坐标；}
  \qpart{（2）求四面体 \(ACB_1D_1\) 六条棱的长，并判断该四面体的形状；}
  \qpart{（3）点 \(G\) 在 \(\triangle CB_1D_1\) 内，且满足 \(\overrightarrow{GC}+\overrightarrow{GB_1}+\overrightarrow{GD_1}=\overrightarrow{0}\)，求 \(G\) 的坐标；判断向量 \(\overrightarrow{AG}\) 分别与 \(\overrightarrow{CB_1}\)，\(\overrightarrow{CD_1}\) 的关系（用数量积验证），并求 \(\left|\overrightarrow{AG}\right|\)．}
  \ti{16}{\fenzhi{12}直三棱柱 \(ABC-A_1B_1C_1\) 中，\(\angle ABC=90^\circ\)，\(AB=BC=2\)，\(AA_1=3\)．以 \(B\) 为原点，\(BA\)，\(BC\)，\(BB_1\) 所在直线分别为 \(x\) 轴、\(y\) 轴、\(z\) 轴建立空间直角坐标系．点 \(P\) 在侧棱 \(CC_1\) 上，设 \(CP=t\)（\(0\leq t\leq 3\)，\(t\) 为实数）．}
  \qpart{（1）用 \(t\) 表示 \(\left|\overrightarrow{AB}\right|^2\)，\(\left|\overrightarrow{PB}\right|^2\)，\(\left|\overrightarrow{PA}\right|^2\)；}
  \qpart{（2）判断 \(\triangle ABP\) 能否为直角三角形与等腰三角形：判断 \(\angle ABP\) 与 \(t\) 的关系并指出直角顶点；求所有使 \(\triangle ABP\) 为等腰三角形的 \(t\) 值；}
  \qpart{（3）记 \(\angle APB=\theta\)，求 \(\tan\theta\) 的取值范围，并求 \(\theta\) 的最大值．}
  \dabiao
}
\end{document}
